\chapter{分布、Green 算子与反问题}
\label{chap:math-distributions}

上一章把点源写成 \(L_xG(x,y)=\delta_y\)。它有效地编码了导数跳跃，但 \(\delta_y\) 不是普通函数：一个普通可积函数若只在单点非零，其积分为零，不能表示单位强度点源。要让这个等式有精确意义，应从“源怎样作用在平滑探针上”定义它。

\section{测试函数与分布}

\begin{definition}[测试函数和收敛]
对开集 \(\Omega\subset\mathbb R^n\)，\(\mathcal D(\Omega)=C_c^\infty(\Omega)\) 是所有支集紧且包含在 \(\Omega\) 内的无限次可微函数的空间。\(\varphi_j\to0\) 于 \(\mathcal D(\Omega)\)，意为所有支集包含在同一个紧集 \(K\Subset\Omega\) 内，且每个多重指标 \(\alpha\) 都满足 \(\sup_K|\partial^\alpha\varphi_j|\to0\)。
\end{definition}

紧支集让分部积分不产生区域边界项；共同的 \(K\) 则保证“探针趋于零”不会同时跑到无穷远处。这比只要求逐点趋零强。例如 \(\varphi_j(x)=\varphi(x-j)\) 在每个固定点最终为零，却因支集逃走而不在 \(\mathcal D(\mathbb R)\) 的上述意义下趋零。

\begin{definition}[分布]
\(\mathcal D(\Omega)\) 上连续的复线性泛函称为分布。局部可积函数 \(f\) 定义 \(\langle T_f,\varphi\rangle=\int_\Omega f(x)\varphi(x)\,\dd x\)；点源定义 \(\langle\delta_y,\varphi\rangle=\varphi(y)\)。两者都是分布，但 \(\delta_y\) 不是由普通局部可积函数给出的分布。
\end{definition}

对支集在固定 \(K\) 的 \(\varphi\)，\(|\varphi(y)|\le\sup_K|\varphi|\)，所以点源确实连续。若想把它当作函数，则其在单点外必须几乎处处为零，与 \(\langle\delta_y,\varphi\rangle=\varphi(y)\) 矛盾。点源的“单位强度”是这条作用规则，而不是一个无穷大的函数值。

\section{导数不必在每一点存在}

普通可微函数满足 \(\int f'\varphi=-\int f\varphi'\)，因为测试函数在边界附近为零。右端只需要 \(f\) 局部可积，即使它有跳跃也有意义。这给出导数的新定义，而不是给不可微点强行指定一个数值。

\begin{definition}[分布导数]
对分布 \(T\)，定义 \(\langle\partial_jT,\varphi\rangle=-\langle T,\partial_j\varphi\rangle\)。反复应用得 \(\langle\partial^\alpha T,\varphi\rangle=(-1)^{|\alpha|}\langle T,\partial^\alpha\varphi\rangle\)。
\end{definition}

这仍是连续线性泛函：若 \(\varphi_n\to0\)，它们的导数在共同紧集上也一致趋零，所以 \(\langle T,\partial_j\varphi_n\rangle\to0\)。因此每个分布都可反复求导；这不声称导数是普通函数。

设 \(H(x)=0\) 于 \(x<0\)，\(H(x)=1\) 于 \(x>0\)。按定义，
\[
\langle H',\varphi\rangle=-\int_0^\infty\varphi'(x)\,\dd x
=-\bigl[\varphi(x)\bigr]_{0}^{\infty}=\varphi(0)
=\langle\delta_0,\varphi\rangle.
\]
故 \(H'=\delta_0\)。更一般地，若 \(f\) 在原点两侧各有普通导数和单侧极限，分段积分分部给
\[
\partial_xT_f=T_{f'_{\mathrm{pw}}}+[f]_0\delta_0,
\qquad [f]_0=f(0^+)-f(0^-).
\]
跳跃的正负号由两侧积分的端点项决定，不能凭图像猜。取 \(f=H\) 时 \([f]_0=1\)，与直接计算一致。

点源也可由普通函数序列逼近，但这种逼近不是它的定义。令 \(\rho_n(x)=n\mathbf1_{[0,1/n]}(x)\)。对每个测试函数 \(\varphi\)，换元 \(s=nx\) 得
\[
\langle T_{\rho_n},\varphi\rangle
=\int_0^1\varphi(s/n)\,\dd s\longrightarrow\varphi(0)
=\langle\delta_0,\varphi\rangle.
\]
极限使用的是 \(\varphi\) 在零点连续，而 \(\rho_n\) 在零点的函数值是否选为 \(n\) 不影响结果。故“高度趋无穷”只能作某些近似列的图像说明；不同归一化可能给零、两倍点源或根本没有分布极限。

\section{Fourier 变换为何要扩大测试函数空间}

紧支测试函数的 Fourier 变换通常不再紧支。为使变换前后仍在同一类探针中，取 Schwartz 空间 \(\mathcal S(\mathbb R^n)\)：其中 \(\varphi\) 光滑，且 \(x^\beta\partial^\alpha\varphi(x)\) 对每组多重指标 \(\alpha,\beta\) 都有界。它的连续对偶称缓增分布。紧支分布和至多多项式增长的函数都在其中，而 \(\ee^{x^2}\) 不属于这一类。

采用 \(\widehat\varphi(\xi)=\int_{\mathbb R^n}\ee^{-\ii x\cdot\xi}\varphi(x)\,\dd^nx\) 的约定，其中 \(x\) 与 \(\xi\) 分别是位置和逆长度波数。反变换必须带 \((2\pi)^{-n}\)：
\begin{equation}
 \varphi(x)=\frac1{(2\pi)^n}\int_{\mathbb R^n}
 \ee^{\ii x\cdot\xi}\widehat\varphi(\xi)\,\dd^n\xi.
 \label{eq:math-fourier-convention}
\end{equation}
先在一维检验这些积分确实有意义。\(\partial_\xi^m\widehat\varphi\) 是 \((-\ii x)^m\varphi\) 的 Fourier 变换；再对 \(x\) 分部积分，可把任意 \(\xi^j\) 移到 \(x\) 导数上。所得被积函数仍可积，所以 \(\widehat\varphi\) 的各阶导数比任意 \(|\xi|^{-j}\) 衰减，\(\widehat\varphi\in\mathcal S\)。要证明反变换而不是把它当形式 delta 恒等式，在 \(\xi\) 积分中临时乘 \(\ee^{-\varepsilon\xi^2}\)，\(\varepsilon>0\)。交换绝对收敛的积分并计算高斯积分，左侧成为
\[
 \frac1{2\pi}\int\ee^{\ii x\xi-\varepsilon\xi^2}\widehat\varphi(\xi)\,\dd\xi
 =\int_{\mathbb R}\frac{\ee^{-(x-y)^2/(4\varepsilon)}}{\sqrt{4\pi\varepsilon}}
 \varphi(y)\,\dd y.
\]
右端高斯核非负、积分为一且宽度趋零，故趋于 \(\varphi(x)\)；左端因 \(\widehat\varphi\in L^1\) 可用有界收敛去掉调节因子。这证明一维反变换；对每个坐标重复同一步便得\cref{eq:math-fourier-convention} 的 \(n\) 维形式。特别地 \(\mathcal F:\mathcal S\to\mathcal S\) 是双射。

对 \(f,g\in\mathcal S(\mathbb R^n)\)，把 \(g\) 的反变换代入 \(\int\overline f g\) 并交换积分，得到
\begin{equation}
 \int_{\mathbb R^n}\overline{f(x)}g(x)\,\dd^nx
 =\frac1{(2\pi)^n}\int_{\mathbb R^n}
 \overline{\widehat f(\xi)}\widehat g(\xi)\,\dd^n\xi.
 \label{eq:math-fourier-plancherel}
\end{equation}
这就是 Parseval 等式的连续变量版本。\(\mathcal S\) 在 \(L^2\) 中稠密，等式因而把 \((2\pi)^{-n/2}\mathcal F\) 唯一延拓为 \(L^2\) 上的等距满射。对一般 \(L^2\) 波函数，Fourier 变换可以按这个范数极限理解，并不要求它的积分在每一点都绝对收敛。

对缓增分布定义 \(\langle\widehat T,\varphi\rangle=\langle T,\widehat\varphi\rangle\)；在普通函数可积时，交换两次积分便恢复普通 Fourier 积分。对任意 \(\varphi\in\mathcal S\)，由 \(\partial_x\widehat\varphi(x)=\widehat{(-\ii\xi\varphi)}(x)\) 得
\[
\begin{aligned}
\langle\widehat{\partial_xT},\varphi\rangle
&=\langle\partial_xT,\widehat\varphi\rangle
=-\langle T,\partial_x\widehat\varphi\rangle\\
&=\langle T,\widehat{(\ii\xi\varphi)}\rangle
=\langle\widehat T,\ii\xi\varphi\rangle.
\end{aligned}
\]
这逐步证明了
\[
\mathcal F(\partial_jT)=\ii\xi_j\widehat T.
\]
例如 \(\widehat{\delta_0}=1\)，故 \(\widehat{\delta_0'}=\ii\xi\)。微分化为乘法后，若符号 \(\ii\xi\) 在 \(\xi=0\) 为零，就不能无条件相除；零频处的困难正对应原方程的常数核。

分布导数总能定义，并不意味着分布能任意相乘。\(H\delta_0\) 若照普通乘法写成“\(H(0)\delta_0\)”，会依赖任意指定的 \(H(0)\)；而 \(H\) 作为局部可积函数的分布根本不记录这个点值。因此一般分布乘法需要额外正则性，\(\delta_0^2\) 也没有由上述定义自动给出的意义。在线性方程中，先确认每个乘积和卷积都已定义，再进行形式运算。

\begin{exercise}
对 \(f(x)=H(x)-H(x-1)\)，不借助普通导数在跳点的值，直接求其分布导数。再用 Fourier 变换验证答案中的两个点源系数。
\end{exercise}

\section{基本解：求一次点源，服务所有普通源}

设 \(L\) 是常系数微分算子。若已找到 \(LE=\delta_0\)，则对紧支光滑源 \(f\)，平移和分布导数的定义给 \(L(E*f)=(LE)*f=\delta_0*f=f\)。\(E\) 称基本解。它只规定全空间中一个点源的局部响应；加上任意齐次解仍是基本解，因此边界条件或无穷远行为必须另外选定。

这里的卷积也可以只用“对探针的作用”定义。对 \(f\in C_c^\infty(\mathbb R^n)\)，令 \((E*f)(x)=\langle E_y,f(x-y)\rangle\)。固定 \(x\) 时，\(y\mapsto f(x-y)\) 是合法测试函数。对 \(x\) 求导只会把导数送到测试函数上；若按 \(y\) 分布积分分部，便得 \(\partial_x^\alpha(E*f)=(\partial^\alpha E)*f=E*(\partial^\alpha f)\)。因此 \(L(E*f)=f\) 并非把奇异对象当普通函数硬乘，而是逐个测试函数可核验的等式。两个任意分布未必能卷积：例如常数函数分布 \(1*1\) 的形式积分在全空间发散。若其中一方紧支，另一个分布便只需在有限范围内接受测试，卷积定义可以延拓；这个支集条件是进行后续代数运算的前提。

对三维 Poisson 算子 \(L=-\Delta\)，源点以外的旋转对称齐次解满足 \((r^2E'(r))'=0\)，所以 \(E(r)=a/r+b\)。要求无穷远趋零给 \(b=0\)。把 \(-\Delta E=\delta_0\) 在半径 \(\varepsilon\) 的球内积分，并用散度定理，
\[
1=-\int_{|x|=\varepsilon}\partial_rE\,\dd S
=-4\pi\varepsilon^2\left(-\frac a{\varepsilon^2}\right)=4\pi a.
\]
因此 \(E(x)=1/(4\pi|x|)\)。这个通量计算同时固定了符号、\(4\pi\) 和点源强度；仅写 \(E\propto1/r\) 不能给出可复算的解。

同一个点源方程还说明为什么必须把\emph{无穷远条件}写进 Green 函数的名字。对实数 \(k>0\)，考虑 \((-\Delta-k^2)G=\delta_0\)。源点外的球对称方程是 \((rG)''+k^2(rG)=0\)，故
\[
 G(r)=\frac{A\ee^{\ii kr}+B\ee^{-\ii kr}}r.
\]
在半径趋零的球上积分原方程，\(-k^2G\) 的体积分趋零，而径向通量要求 \(4\pi(A+B)=1\)。这一条件只固定 \(A+B\)，没有选择波向外走还是向内走。若规定向外传播的 Sommerfeld 条件 \(r(\partial_r-\ii k)G(r)\to0\)（\(r\to\infty\)），入射项 \(\ee^{-\ii kr}/r\) 不满足它，故 \(B=0\)，得到
\begin{equation}
 G_+(r)=\frac{\ee^{\ii kr}}{4\pi r}.
 \label{eq:math-outgoing-helmholtz}
\end{equation}
反向条件会选出 \(G_-(r)=\ee^{-\ii kr}/(4\pi r)\)。两者都在局部满足同一分布方程，却对应不同的物理边界条件。由于实 \(k^2\) 落在全空间 \(-\Delta\) 的连续谱上，\(G_+\) 不是 \((-\Delta-k^2)^{-1}\) 在整个 \(L^2\) 上的有界核；它是在辐射条件下定义的响应。若误把它当普通有界逆，便会抹去“选向外波”这一步。

在有边界的区域 \(\Omega\) 内，Green 核满足 \(L_xG(x,y)=\delta_y\) 和指定边界条件。对 Dirichlet 问题，可写 \(G(x,y)=E(x-y)+h_y(x)\)，其中 \(Lh_y=0\)，且边界上的 \(h_y=-E(\cdot-y)\)。第一项负责源点奇性，第二项负责边界。两项缺一不可；只用全空间基本解通常不会满足实验装置的边界。

奇异势提供一个能把“分布”和“预解式”同时算清的例子。在实线上先解 \(L_0=-\partial_x^2+\kappa^2\)，\(\kappa>0\)。在源点外，衰减解与 \(\ee^{-\kappa|x-y|}\) 成比例；其导数在 \(x=y\) 的跳跃为 \(-2\kappa\)，要求 \(-[\partial_xG_0]_{y}=1\) 给
\[
G_0(x,y)=\frac{\ee^{-\kappa|x-y|}}{2\kappa}.
\]
现在加入强度 \(g\in\mathbb R\) 的点势，方程 \((-\partial_x^2+\kappa^2)u+g\delta_0u=f\) 只对在零点连续的 \(u\) 使用 \(\delta_0u=u(0)\delta_0\)。由已知的 \(G_0\) 求解得到 \(u=G_0f-gG_0(\cdot,0)u(0)\)。取 \(x=0\) 解一个标量方程：
\[
u(0)=\frac{(G_0f)(0)}{1+g/(2\kappa)}.
\]
因此只要 \(2\kappa+g\ne0\)，新 Green 核为
\[
G_g(x,y)=G_0(x,y)-\frac{g}{1+g/(2\kappa)}G_0(x,0)G_0(0,y).
\]
把它在 \(x=0\) 两侧求导，得 \(u'(0^+)-u'(0^-)=g u(0)\)，恰是把原方程跨越零点积分所得的跳跃条件。若 \(g<0\) 且 \(\kappa=-g/2\)，分母为零，出现束缚态 \(\ee^{g|x|/2}\)；预解式极点又一次与本征态对应。

点势的算子本身也必须写出域，才能把这一说法变成谱结论。令 \(H_g u=-u''\) 于 \(x\ne0\)，域由 \(u\in H^1(\mathbb R)\)、两侧各属 \(H^2\)，以及 \(u'(0^+)-u'(0^-)=g u(0)\) 组成。\(H^1\) 使 \(u\) 在零点连续；对分段 \(H^2\) 函数作两次分布求导，\(-u''\) 还含有 \(-[u']_0\delta_0\)，恰与 \(g u(0)\delta_0\) 相消。这个域不是任意补上的：形式
\[
 q_g[u,v]=\int_{\mathbb R}\overline{u'}v'\,\dd x
   +g\overline{u(0)}v(0),\qquad D(q_g)=H^1(\mathbb R)
\]
在实 \(g\) 时对称，点值项由一维迹估计相对于 \(\|u'\|_2^2+\|u\|_2^2\) 控制；前述表示定理给出唯一自伴实现，其域条件正是刚算出的导数跳跃。当 \(g<0\)，归一化态 \(u_b(x)=\sqrt{-g/2}\,\ee^{g|x|/2}\) 满足 \(H_gu_b=-g^2u_b/4\)，直接积分 \(\int_{\mathbb R}|u_b|^2=1\) 可核验系数。注意这里 \(H_g\) 的负本征值是 \(-g^2/4\)；上式的预解式分母在 \(\kappa=-g/2\) 为零，正对应 \(H_g+\kappa^2\) 失去可逆性。

\begin{exercise}
验证二维函数 \(E(x)=-(2\pi)^{-1}\log|x|\) 在原点外调和，并在半径 \(\varepsilon\) 的圆上计算 \(-\partial_rE\) 的总通量，说明为什么它满足 \(-\Delta E=\delta_0\)。
\end{exercise}


\section{基本解什么时候能与源卷积}

在 \(\mathbb R^n\) 上，常系数微分算子 \(L\) 的基本解 \(E\) 满足 \(LE=\delta_0\)。对足够好的源 \(f\)，平移不变性给 \(u=E*f\)，且 \(Lu=(LE)*f=f\)。若 \(f\) 是光滑紧支函数，卷积可用分布作用 \((E*f)(x)=\langle E_y,f(x-y)\rangle\) 定义。对两个任意分布则未必能卷积；若至少一个紧支，定义可以通过测试函数的双重作用延拓。边界 Green 核 \(G(x,y)\) 不具平移不变性，故通常写 \(u(x)=\int G(x,y)f(y)\,\dd y\)，而非 \(G*f\)。

\section{Green 算子遇到零模时怎样求解}

若 \(L_0u=f\) 的边界 Green 算子为 \(G_0\)，加上势项后 \((L_0+V)u=f\) 等价于
\[
u+G_0Vu=G_0f.
\]
这是第二类积分方程。若 \(G_0V\) 紧，它看起来像有限维的“单位矩阵加一个可控修正”，但可能出现非零齐次解；这时不能只宣布逆不存在，还要问哪些源仍能解。

\begin{theorem}[Fredholm 择一原理]
若 \(K:\mathcal H\to\mathcal H\) 紧，则 \(T=I-K\) 的核有限维，像闭，且 \(\dim\ker T=\dim\ker T^*\)。方程 \(Tu=f\) 可解当且仅当 \(f\perp\ker T^*\)。特别地，若 \(\ker T=\{0\}\)，则对每个 \(f\) 有唯一解。
\end{theorem}

先说明紧性在证明中做了什么。若 \(\ker T\) 无限维，可在其中选正交单位列 \(e_n\)。由 \(Te_n=0\) 知 \(Ke_n=e_n\)，但 \((Ke_n)\) 没有收敛子列，矛盾。故核有限维。若像不闭，可在 \((\ker T)^\perp\) 里选 \(\|x_n\|=1\)、\(Tx_n\to0\)。由紧性取子列使 \(Kx_n\) 收敛；于是 \(x_n=Tx_n+Kx_n\) 也收敛到一个核内单位向量，却仍属于核的正交补，矛盾。因此像闭。

对任何有界算子，\((\operatorname{Ran}T)^\perp=\ker T^*\)：\(f\) 与所有 \(Tu\) 正交，恰等于 \(T^*f=0\)。像既闭，便有 \(\operatorname{Ran}T=(\ker T^*)^\perp\)，这证明了可解性的充要条件。还差核与余核维数相等；紧性允许把它严格化为有限维线性代数。

先证有限秩逼近的来源。对 \(K\) 作用于单位球所得集合，紧性给出有限个半径 \(\varepsilon\) 的球覆盖，设球心张成有限维空间 \(M\)，令 \(P_M\) 为其正交投影。每个 \(Kx\) 与 \(M\) 距离小于 \(\varepsilon\)，故 \(F=P_MK\) 有限秩且 \(\|K-F\|\le\varepsilon\)。取 \(\varepsilon<1\)，写 \(R=K-F\)。Neumann 级数 \(\sum_{j\ge0}R^j\) 在算子范数下收敛，给出 \((I-R)^{-1}\)；于是
\[
T=I-K=(I-R)(I-G),\qquad G=(I-R)^{-1}F.
\]
\(G\) 仍有限秩，前一因子可逆，所以 \(T\) 与 \(I-G\) 的核同维、像的余维也相同。令 \(M_G=\operatorname{Ran}G\) 并正交分解 \(\mathcal H=M_G\oplus M_G^\perp\)，则 \(I-G\) 的块矩阵为
\[
\begin{pmatrix}I_{M_G}-G_{11}&-G_{12}\\0&I_{M_G^\perp}\end{pmatrix}.
\]
其核维数等于有限矩阵 \(I_{M_G}-G_{11}\) 的核维数；像的余维也等于该有限矩阵的像在 \(M_G\) 中的余维。有限维秩零度定理使两者相等。因此 \(\dim\ker T=\dim\ker T^*\)，也证明了定理中“单射便满射”的结论。

取 \(\mathcal H=\ell^2\)，\(Kx=\langle e_1,x\rangle e_1\)。则 \(Tu=(0,u_2,u_3,\ldots)\)，齐次解是 \(e_1\) 的任意倍；\(Tu=f\) 当且仅当 \(f_1=0\)，其余分量由 \(u_j=f_j\) 给出。这里条件 \(f\perp\ker T^*\) 就是熟悉的“右端与左零空间正交”。在 Neumann Laplacian 中，它变为 \(\int_0^1f(x)\,\dd x=0\)。

紧性不能删。令 \(S(x_1,x_2,\ldots)=(0,x_1,x_2,\ldots)\) 为 \(\ell^2\) 上的右移算子，它有界但不紧。\((I-S)x=0\) 逐分量迫使 \(x=0\)，可 \((I-S)x=e_1\) 要求每个 \(x_n=1\)，所得序列不在 \(\ell^2\)。所以 \(I-S\) 单射而不满射；有限维直觉恰在缺少紧性的地方失效。

\begin{exercise}
令 \(Kx=\lambda\langle e_1,x\rangle e_1\)。分别讨论 \(\lambda=1\) 与 \(\lambda\ne1\) 时 \((I-K)u=f\) 的存在性和唯一性。
\end{exercise}

\section{逆问题：数据误差怎样被小奇异值放大}

正问题 \(y=Kx\) 若 \(K\) 有界，则输入误差至多被 \(\|K\|\) 放大。反过来由 \(y\) 求 \(x\) 却可能极不稳定。取 \(\ell^2\) 上的紧算子 \((Kx)_n=x_n/n\)。对无噪数据 \(y\) 若确实来自某个 \(x\)，形式逆是 \(x_n=ny_n\)。把仅第 \(N\) 个分量为 \(1/N\) 的噪声加到零数据上，数据范数趋零，而反演结果范数恒为一。故逆映射不连续；这与计算机精度无关。

\begin{definition}[奇异系统]
对紧算子 \(K:X\to Y\)，若正数 \(s_n\) 与单位正交向量 \(v_n\in X\)、\(u_n\in Y\) 满足 \(Kv_n=s_nu_n\)、\(K^*u_n=s_nv_n\)，则称它们构成奇异系统。\(s_n\) 是奇异值。
\end{definition}

在这些方向上，形式逆把数据系数 \(\langle u_n,y\rangle\) 除以 \(s_n\)。当 \(s_n\to0\) 时，噪声正会在最小奇异值处被放大。截断奇异值法只保留 \(s_n\ge\tau\) 的分量；Tikhonov 方法用上一节的正规方程，给
\[
x_\alpha=\sum_n\frac{s_n}{s_n^2+\alpha}\langle u_n,y\rangle v_n
\quad\text{在 }(\ker K)^\perp\text{ 上}.
\]
系数满足 \(s/(s^2+\alpha)\le1/(2\sqrt\alpha)\)，所以噪声 \(\|y^\delta-y\|\le\delta\) 导致的解误差至多为 \(\delta/(2\sqrt\alpha)\)。这是稳定性的定量结论，但 \(\alpha\) 越大，对真实信号的抑制也越强。

为什么紧算子能按这些方向拆开，也值得从已经证明的谱定理核对一次。\(K^*K:X\to X\) 是紧、自伴且非负的：紧性由有界算子与紧算子的复合得到，\(\langle x,K^*Kx\rangle=\|Kx\|^2\ge0\)。在 \((\ker K)^\perp\) 上，紧自伴谱定理给标准正交本征向量 \(v_n\) 和正本征值 \(s_n^2\)，并且 \(s_n\downarrow0\)（有限秩时序列在有限项后结束）。设 \(u_n=Kv_n/s_n\)，则
\[
 \langle u_m,u_n\rangle
 =\frac{\langle v_m,K^*Kv_n\rangle}{s_ms_n}
 =\delta_{mn},
 \qquad K^*u_n=s_nv_n.
\]
若 \(y\perp u_n\) 对全部 \(n\) 成立且 \(y\perp\ker K^*\)，则 \(K^*y\) 与全部 \(v_n\) 及 \(\ker K\) 正交，所以 \(K^*y=0\)，进一步得到 \(y=0\)。因此 \((u_n)\) 张成 \(\overline{\operatorname{Ran}K}\)，\(y\) 在这部分之外的分量根本不是正问题可能产生的数据。

这给出一个可检验的可解性判据，而不只是“不稳定”的形容。若 \(y\in Y\)，方程 \(Kx=y\) 有解的充要条件是
\begin{equation}
 y\perp\ker K^*,\qquad
 \sum_n\frac{|\langle u_n,y\rangle|^2}{s_n^2}<\infty.
 \label{eq:math-picard-condition}
\end{equation}
必要性由 \(x=\sum_n\langle v_n,x\rangle v_n+x_0\)、\(x_0\in\ker K\) 代入 \(Kx\) 与 Parseval 立得。反过来，若两条件成立，令 \(x=\sum_n s_n^{-1}\langle u_n,y\rangle v_n\)。第二条件保证这列向量在 \(X\) 中收敛，\(K\) 有界允许逐项作用；第一条件和 \((u_n)\) 在 \(\overline{\operatorname{Ran}K}\) 的完备性再保证 \(Kx=y\)。这个结论也说明\emph{存在性与稳定性是两件事}：即使精确数据满足平方和条件，任意小的高序号噪声也可能使该和发散。

例如 \(s_n=1/n\)，取 \(y_n=1/n^2\)。此时 \(y\in\ell^2\)，且逆解 \(x_n=1/n\) 仍在 \(\ell^2\)，所以方程确实可解。若改取 \(y_n=1/n^{3/2}\)，数据仍平方可和，但形式逆为 \(x_n=1/\sqrt n\)，其平方和发散；数据处在 \(\overline{\operatorname{Ran}K}\) 却不在实际的像中。后面加正则化时，即使 \(y\) 不在像中，\(x_\alpha\) 仍由正规方程唯一给出；它求的是带偏差的最佳折中，不是原方程凭空出现了精确解。

对角例 \(s_n=1/n\) 的每个分量可直接写成
\[
(x_\alpha)_n=\frac{n}{1+\alpha n^2}\,y_n^\delta.
\]
低序号近似乘以 \(n\)，高序号 \(n\gg\alpha^{-1/2}\) 则被压低。这让正则化参数具有明确的分辨率含义。若把 \(\alpha\) 无条件送到零，噪声放大又会回来；极限必须与噪声水平一起讨论。

把对角例换成真正的模糊成像也会出现同一结构。在周期区间 \([0,2\pi]\) 上，取 Fourier 基 \(e_n(x)=(2\pi)^{-1/2}\ee^{\ii nx}\)。若仪器的卷积核把第 \(n\) 个频率乘以 \(s_n=\ee^{-\sigma^2n^2/2}\)，则 \(Ke_n=s_ne_n\)。它平滑高频，因此是稳定的正问题；形式逆却乘 \(\ee^{\sigma^2n^2/2}\)，比前面 \(n\) 的放大还快。对带噪观测 \(y^\delta\)，Tikhonov 给每个 Fourier 模的重建系数
\[
\langle e_n,x_\alpha\rangle
=\frac{s_n}{s_n^2+\alpha}\langle e_n,y^\delta\rangle.
\]
当 \(s_n^2\ll\alpha\) 时系数被压低；这就是空间分辨率损失的数学含义。任意算法都无法从噪声中可靠找回这些方向，除非加入额外先验。

参数怎样随噪声减小，可在一个明确的光滑先验下证明。假设真实 \(x^\dagger=(K^*K)^\nu w\)，其中 \(0<\nu\le1\)、\(\|w\|\le M\)，且 \(x^\dagger\perp\ker K\)。无噪 Tikhonov 系数与真值之差的滤波因子为 \(\alpha/(s_n^2+\alpha)\)；把 \(x^\dagger\) 的系数 \(s_n^{2\nu}\langle v_n,w\rangle\) 代入，并令 \(t=s_n^2/\alpha\)，得到偏差界
\[
\|x_\alpha(y)-x^\dagger\|
\le M\alpha^\nu\sup_{t\ge0}\frac{t^\nu}{1+t}
\le M\alpha^\nu.
\]
加上前面的噪声界，\(\|x_\alpha(y^\delta)-x^\dagger\|\le M\alpha^\nu+\delta/(2\sqrt\alpha)\)。平衡两项可取 \(\alpha\) 与 \(\delta^{2/(2\nu+1)}\) 同阶，误差便与 \(\delta^{2\nu/(2\nu+1)}\) 同阶。没有这个光滑先验就不能宣称这一速率；若真值含 \(\ker K\) 分量，数据根本不含那部分信息。

\begin{exercise}
取真实信号 \(x=e_N\)，则无噪数据 \(y=e_N/N\)。计算上式给出的第 \(N\) 个分量及偏差。再把同方向的噪声 \(\eta=e_N/N\) 加到零数据，比较偏差与噪声放大。
\end{exercise}


\section{正则化方程为什么有唯一解}

\begin{proposition}[Tikhonov 正规方程]
设 \(K:X\to Y\) 为 Hilbert 空间间的有界线性算子，\(y\in Y\)，\(\alpha>0\)。泛函 \(J_\alpha(x)=\|Kx-y\|_Y^2+\alpha\|x\|_X^2\) 有唯一极小点 \(x_\alpha\)，且
\[
(K^*K+\alpha I)x_\alpha=K^*y.
\]
\end{proposition}
\begin{proof}
写 \(B=K^*K+\alpha I\)。它有界、自伴，且 \(\langle x,Bx\rangle=\|Kx\|^2+\alpha\|x\|^2\ge\alpha\|x\|^2\)。Cauchy--Schwarz 推出 \(\|Bx\|\ge\alpha\|x\|\)，故 \(B\) 的像闭且核为零。又 \(\operatorname{Ran}(B)^\perp=\ker B^*=\ker B=\{0\}\)，所以像也稠密，只能等于 \(X\)。正规方程因而有唯一解。令该解为 \(x_\alpha\)，对任意 \(h\in X\) 展开平方并使用 \(Bx_\alpha=K^*y\)，得到
\[
J_\alpha(x_\alpha+h)-J_\alpha(x_\alpha)
=\|Kh\|^2+\alpha\|h\|^2.
\]
右端仅在 \(h=0\) 时为零，故解是唯一极小点。
\end{proof}

正规方程使不稳定的形式逆变成稳定的有界求解，但 \(\alpha\) 也引入偏差。它与噪声水平的关系要在具体奇异值算例里看清。
